\section{Extending the comparison and descending sections}
\label{descent:section}

Two analytic facts complete the passage from the auxiliary parameter
space to the original base model.  The first extends a specified
isomorphism of volume lines; the second descends global generation.

\begin{lemma}
\label{descent:unit}
Let \(P\) be a complex manifold, \(D\) a simple normal crossing
divisor, and \(P^\circ=P\setminus D\).
Let \(E_1,E_2\) be holomorphic line bundles on \(P\), equipped on
\(P^\circ\) with Hermitian norms, and let
\[
 \phi:E_1|_{P^\circ}\xrightarrow{\ \sim\ }E_2|_{P^\circ}
\]
be a specified holomorphic isomorphism.
Assume the following near the smooth locus of each component of
\(D\).
In a coordinate polydisk \((t,z)\) with \(D=(t=0)\), choose
nonvanishing holomorphic frames \(e_1,e_2\).
For a dense set of centers \(z\), their norms and dual norms along
the parallel punctured disks satisfy bounds by powers of
\(1+|\log|t||\).  The constants and powers may depend on the disk.
Assume also that \(\phi\) compares the two norms by factors bounded
above and below on each such disk.
It is enough for these bounds to hold after a finite power
substitution on each disk.
Then \(\phi\) extends uniquely to a holomorphic isomorphism
\(E_1\simeq E_2\) on \(P\).
\end{lemma}

\begin{proof}
Write \(\phi(e_1)=g e_2\), where \(g\) is holomorphic and
nowhere zero on the punctured polydisk.
The bounds on the frames and their duals give, on every disk under
consideration,
\begin{equation}
\label{descent:scalar-bounds}
 |g(t,z)|+|g(t,z)|^{-1}
 \leq C_z\bigl(1+|\log|t||\bigr)^{b_z}
\end{equation}
for suitable positive constants.
If the estimate is initially given after \(t=u^N\), it implies
an estimate of the same form downstairs, since every nonzero \(t\)
has an \(N\)-th root and
\(|\log|u||=|\log|t||/N\).

Fix one such \(z\).
For \(j\geq1\), the coefficient of \(t^{-j}\) in the Laurent
expansion of \(g(t,z)\) is
\[
 a_{-j}(z)=\frac{1}{2\pi i}
          \int_{|t|=\varepsilon}g(t,z)t^{j-1}\,dt.
\]
It is independent of \(\varepsilon\), and
\eqref{descent:scalar-bounds} yields
\[
 |a_{-j}(z)|
 \leq C_z\varepsilon^j
        \bigl(1+|\log\varepsilon|\bigr)^{b_z}
 \longrightarrow0.
\]
Thus every negative Laurent coefficient vanishes.
The same argument applies to \(g^{-1}\), so both functions extend
holomorphically on this disk, and their extensions multiply to one.

For a fixed integration radius, \(a_{-j}(z)\) is holomorphic in
the remaining coordinates.  It vanishes on the dense set of
centers from the hypothesis and therefore vanishes identically.
The Laurent expansion in the polydisk consequently has no negative
powers.  Equivalently, Cauchy's integral formula constructs a
jointly holomorphic extension across \(t=0\).
Applying this also to \(g^{-1}\) gives a holomorphic unit throughout
the polydisk.  This reasoning does not require uniform logarithmic
constants as the center varies.

We have extended \(\phi\) and its inverse across each component of
\(D\) away from its intersections with the others.
The remaining subset has complex codimension at least two.
In local line-bundle frames, Hartogs' extension theorem extends the
coefficient functions of both maps across that subset.
Their compositions remain the identity by analytic continuation.
Local extensions agree because they agree on the dense open
\(P^\circ\), proving existence and uniqueness globally.
\end{proof}

\begin{remark}
\label{descent:tangential-disks}
Local power charts also handle disks tangent to the boundary.
Indeed, write such a chart as
\[
 (u_1,\ldots,u_\ell,z)\longmapsto
 (u_1^{N_1},\ldots,u_\ell^{N_\ell},z).
\]
For a holomorphic disk \(\gamma\) whose punctured disk avoids the
boundary, each boundary coordinate has the form
\[
 t_i\circ\gamma(t)=t^{b_i}a_i(t),
 \qquad b_i\geq0,\quad a_i(0)\ne0.
\]
After a substitution \(t=v^N\), with every \(N_i\) dividing \(N\),
the functions \(a_i(v^N)\) have holomorphic \(N_i\)-th roots on a
smaller disk, so \(\gamma\) lifts to the power chart.
Thus the disk bounds for generators supplied by
Theorem~\ref{orders:trace} and the factor-extension results can be
used after arbitrary auxiliary pullbacks, including at tangential
disks.  Lemma~\ref{descent:unit} itself needs only the general
parallel transverse disks.
\end{remark}

We next give the descent statement in a form that does not require
either a generically finite auxiliary map or a finite flat
intermediate map.

\begin{lemma}
\label{descent:semiample}
Let \(r:P\to S\) be a proper surjective holomorphic map of
connected compact complex manifolds, and let \(A\) be a
holomorphic line bundle on \(S\).
If \(r^*A\) is semiample, then \(A\) is semiample.
More precisely, let
\[
 P\xrightarrow{h}T\xrightarrow{q}S
\]
be the Stein factorization, and let \(d\) be the generic degree of
\(q\).
If \(r^*(A^{\otimes n})\) is globally generated, then
\(A^{\otimes nd}\) is globally generated.
The same descent conclusion holds for rational holomorphic line
bundles after clearing denominators.
\end{lemma}

\begin{proof}
The intermediate space \(T\) is normal and irreducible,
\(q\) is finite surjective, and \(h_*\cO_P=\cO_T\).
The projection formula gives an identification of section spaces
\begin{equation}
\label{descent:stein-sections}
 H^0\bigl(P,r^*(A^{\otimes n})\bigr)
 \simeq
 H^0\bigl(T,q^*(A^{\otimes n})\bigr).
\end{equation}
For \(t\in T\), choose \(p\in h^{-1}(t)\).
A section on the left that is nonzero at \(p\) descends to a
section on the right that is nonzero at \(t\).
Thus \(q^*(A^{\otimes n})\) is globally generated.

We recall directly the analytic norm needed for \(q\), without
assuming flatness.
Outside a proper analytic subset of \(S\), the map \(q\) is a
covering with \(d\) distinct sheets.
If \(b\) is a holomorphic function on \(q^{-1}(V)\), its product
over these sheets is a single-valued holomorphic function on the
covering locus in \(V\):
\[
 \Nm_q(b)(s)=\prod_{t\in q^{-1}(s)} b(t).
\]
This function is locally bounded near the omitted analytic subset.
To see this near \(s_0\), choose neighborhoods of the finitely many
points of \(q^{-1}(s_0)\) on which \(b\) is bounded.
Properness allows us to shrink \(V\) so that \(q^{-1}(V)\) is
contained in their union.
The product is then bounded by the \(d\)-th power of a common
bound.  Since \(S\) is normal, the removable singularities theorem
extends it uniquely to a holomorphic function on \(V\).
The extended norm is multiplicative and satisfies
\[
 \Nm_q(q^*c)=c^d.
\]
Both identities hold on the covering locus and hence everywhere.

If \(b\) is nonzero at every point above \(s_0\), the same
neighborhoods can be chosen so that \(b\) and \(b^{-1}\) are
bounded.  The preceding construction then extends both
\(\Nm_q(b)\) and \(\Nm_q(b^{-1})\), whose product is one.
In particular, \(\Nm_q(b)(s_0)\ne0\).
This argument uses local sheets and normality, rather than a
determinant of a locally free pushforward.

For any section
\[
 v\in H^0\bigl(T,q^*(A^{\otimes n})\bigr),
\]
apply the function norm to its coefficients in local frames of
\(A^{\otimes n}\).
The identity \(\Nm_q(q^*c)=c^d\) shows that these coefficients glue
to a section
\[
 \Nm_q(v)\in H^0\bigl(S,A^{\otimes nd}\bigr).
\]
Fix \(s_0\in S\).
For each point \(t\) of the finite set \(q^{-1}(s_0)\), sections
vanishing at \(t\) form a proper linear hyperplane in the section
space in \eqref{descent:stein-sections}.
Over \(\C\), finitely many proper hyperplanes cannot cover a
vector space.  We may therefore choose \(v\) nonzero at every
point of \(q^{-1}(s_0)\).
Its norm is nonzero at \(s_0\).
The exponent \(nd\) is independent of \(s_0\), so these norm
sections prove global generation of \(A^{\otimes nd}\).

Finally, for a rational line bundle, first choose an actual line
bundle representing a positive integral multiple.
Any torsion ambiguity in such a representative disappears after
a further positive tensor power.  The preceding argument then
applies to that actual line bundle.
\end{proof}

\begin{remark}
No Kähler assumption on the auxiliary manifold \(P\) is used in
Lemma~\ref{descent:semiample}.  Positive-dimensional fibers are
handled by \eqref{descent:stein-sections}; branching and possible
nonflatness of the finite map are handled by the bounded norm.
In particular, this descent gives global generation on the whole
compact base, with a single exponent.
\end{remark}

